\chapter{Airy Disk}

\section*{Introduction}

The Airy disk is the diffraction pattern produced by a circular aperture when illuminated by a point source. Named after George Biddell Airy, who derived its mathematical description in 1835, this pattern is the fundamental resolution limit of all optical instruments---telescopes, microscopes, cameras, and the human eye. The bright central spot, surrounded by concentric rings of decreasing intensity, sets the minimum angular size of two point sources that can be distinguished as separate. Every photograph ever taken, every image ever formed by a lens, is ultimately limited by the Airy disk.

When light passes through a circular opening (the aperture of a telescope, the pupil of an eye), it diffracts---the wavefronts spread out and interfere with each other. For a perfect circular aperture, the resulting intensity pattern on a distant screen consists of a bright central disk (the Airy disk) surrounded by concentric bright and dark rings. The central disk contains about 84\% of the total light energy. Its angular radius is set by the ratio of the wavelength to the aperture diameter: the smaller the aperture (relative to the wavelength), the larger the diffraction pattern and the worse the resolution. This is why large telescopes see finer details than small ones, and why short-wavelength (blue/UV) light can resolve smaller features than long-wavelength (red/IR) light.
\section*{Fundamental Equation Used}

\begin{equation}
\theta = 1.22\,\frac{\lambda}{D}
\end{equation}
\section*{Assumptions}

\textbf{Assumptions:} The aperture is perfectly circular. The incident light is monochromatic and spatially coherent (a point source at infinity, or a laser beam). The optical system is aberration-free (the Airy disk is the diffraction limit, assuming no lens imperfections). The observation is in the far field (Fraunhofer regime), i.e., the screen distance satisfies $L \gg D^2/\lambda$.


\textbf{Limitations:} Real optical systems have aberrations (spherical, chromatic, coma) that may make the actual spot larger than the Airy disk. Polychromatic light (white light) produces overlapping Airy patterns for each wavelength, blurring the rings. Off-axis points produce distorted (coma) patterns. The Rayleigh criterion ($\theta = 1.22\lambda/D$) is a convention, not a fundamental limit; super-resolution techniques can beat it under specific conditions.


\section*{Mathematical Derivation}

\textbf{Step 1: Define the aperture function.}

Consider a circular aperture of diameter $D$ (radius $R = D/2$) in the $xy$-plane, illuminated by a normally incident plane wave of wavelength $\lambda$. The aperture function is:
\begin{equation}
P(x,y) = \begin{cases} 1 & \text{if } x^2 + y^2 \leq R^2 \\ 0 & \text{otherwise} \end{cases}
\end{equation}

\textbf{Step 2: Write the Fraunhofer diffraction integral.}

In the far field (Fraunhofer regime), the diffracted electric field at angular position $(\theta_x, \theta_y)$ is proportional to the 2D Fourier transform of the aperture function:
\begin{equation}
E(\theta_x, \theta_y) \propto \iint_{\text{aperture}} P(x,y)\,e^{-i k(\theta_x x + \theta_y y)}\,dx\,dy
\end{equation}
where $k = 2\pi/\lambda$.

\textbf{Step 3: Convert to polar coordinates in the aperture plane.}

Exploiting the circular symmetry of the aperture, switch to polar coordinates $(r,\varphi)$ in the aperture and $(\theta,\phi)$ in the observation direction:
\begin{equation}
E(\theta) \propto \int_0^R \int_0^{2\pi} e^{-ikr\theta\cos(\varphi-\phi)}\,r\,dr\,d\varphi
\end{equation}
The $\varphi$-integral gives a Bessel function:
\begin{equation}
\int_0^{2\pi} e^{-ikr\theta\cos(\varphi-\phi)}\,d\varphi = 2\pi J_0(kr\theta)
\end{equation}

\textbf{Step 4: Perform the radial integration.}

The radial integral becomes:
\begin{equation}
E(\theta) \propto 2\pi\int_0^R J_0(kr\theta)\,r\,dr
\end{equation}
Using the identity $\int_0^a J_0(\alpha r)\,r\,dr = \frac{a}{\alpha}J_1(\alpha a)$ with $\alpha = k\theta$:
\begin{equation}
E(\theta) \propto 2\pi\frac{R}{k\theta}J_1(kR\theta)
\end{equation}

\textbf{Step 5: Normalize and form the intensity.}

Define the dimensionless variable $x = kR\theta = \pi D\sin\theta/\lambda$. The normalized electric field is:
\begin{equation}
E(\theta) \propto \frac{2J_1(x)}{x}
\end{equation}
The intensity is the squared modulus:
\begin{equation}
I(\theta) = I_0\left[\frac{2J_1(x)}{x}\right]^2
\end{equation}
where $I_0$ is the intensity at the center ($\theta = 0$). At $x = 0$, $J_1(x)/x \to 1/2$, confirming $I(0) = I_0$.

\textbf{Step 6: Find the first dark ring.}

The first zero of $J_1(x)$ occurs at $x_1 = 3.8317$. Setting $x = \pi D\sin\theta/\lambda = x_1$:
\begin{equation}
\sin\theta_{\text{first zero}} = \frac{3.8317}{\pi}\frac{\lambda}{D} = 1.220\,\frac{\lambda}{D}
\end{equation}
For small angles ($\sin\theta \approx \theta$):
\begin{equation}
\theta_{\text{Airy}} = 1.22\,\frac{\lambda}{D}
\end{equation}
This is the Rayleigh criterion for the angular resolution of a circular aperture.

\textbf{Step 7: Verify the energy distribution.}

Integrating the intensity over the entire pattern gives the total power. The fraction of power within the first dark ring (the Airy disk) is:
\begin{equation}
\eta = 1 - J_0^2(x_1) - J_1^2(x_1) \approx 0.838
\end{equation}
Thus approximately 84\% of the total light energy is concentrated in the central disk. The first bright ring contains about 1.75\% of the total energy, the second ring 0.42\%, and so on. The subsequent zeros of $J_1(x)$ are at $x = 7.0156, 10.1735, 13.3237, \ldots$, giving rings at angles that are increasingly close to each other.

\section*{Characteristics of the Equation}

The intensity distribution is $I(\theta) = I_0\left[2J_1(x)/x\right]^2$, where $x = \pi D\sin\theta/\lambda$ and $J_1$ is the first-order Bessel function of the first kind. The first zero occurs at $x = 3.8317$, giving $\theta_{\text{first zero}} = 1.22\lambda/D$. The central disk has a diameter (in linear units on a screen at distance $L$) of $d = 2.44\lambda L/D$. The first bright ring has intensity only 1.75\% of the central maximum; the second ring, 0.42\%. The Airy disk is the point spread function (PSF) of a diffraction-limited circular aperture.
\section*{Solution Methods}

The Airy pattern is derived by Fourier transforming the circular aperture function (a circle in the pupil plane transforms to a Bessel function in the focal plane). For a given aperture diameter $D$ and wavelength $\lambda$: compute $\theta_{\text{Airy}} = 1.22\lambda/D$. Convert to linear size on the detector: $d = 2f\theta_{\text{Airy}} = 2.44\lambda f/D$, where $f$ is the focal length. The $f$-number is $N = f/D$, so $d = 2.44\lambda N$. For two point sources, the Rayleigh criterion says they are just resolvable when the center of one Airy disk falls on the first dark ring of the other: $\Delta\theta = 1.22\lambda/D$. The Sparrow criterion (less conservative) sets the limit at $\Delta\theta = 0.94\lambda/D$.
\section*{Verifications}

The Airy pattern has been verified in countless experiments. A laser pointed at a pinhole produces concentric rings whose angular positions match $1.22\lambda/D$. The intensity profile $[2J_1(x)/x]^2$ has been measured with photodetector scans and CCD cameras, confirming the Bessel-function shape to high precision. The Rayleigh criterion is routinely used to calibrate optical systems: two point sources separated by $1.22\lambda/D$ produce a 26\% intensity dip between the two peaks, consistent with the theoretical prediction.
\section*{Applications}

Telescope design (the primary mirror diameter determines the angular resolution: a 10-meter telescope at 500 nm has $\theta_{\text{Airy}} = 0.013''$), microscope resolution (the Abbe limit $d = 0.61\lambda/\text{NA}$ is the Airy-disk result for a lens with numerical aperture NA), camera and sensor design (pixel size should match the Airy disk for optimal sampling---the ``Nyquist criterion''), astronomical imaging (adaptive optics corrects atmospheric turbulence to reach the diffraction limit), laser beam propagation (the divergence of a Gaussian beam is related to the Airy pattern of its truncated aperture), and optical testing (the quality of a lens or mirror is assessed by how closely its point spread function matches the ideal Airy pattern).
\section*{What Richard Feynman Would Have Asked}

\textit{``Why does a circular aperture produce rings? Doesn't the light just spread out in a disk?''}

Feynman would explain this using Huygens' principle: every point on the wavefront passing through the aperture is a source of a secondary wavelet. These wavelets spread out and interfere with each other. At the center of the pattern (straight ahead), all the wavelets arrive in phase and add constructively---this is the bright central maximum. At a slightly off-center point, the wavelets from different parts of the aperture arrive with slightly different phases. At the angle where the path difference between the top and bottom of the aperture is exactly one wavelength, the wavelets from the upper half cancel those from the lower half---this is the first dark ring. Beyond that, there is partial re-cancellation and partial constructive interference, producing the weaker bright rings. The rings are a direct consequence of the circular symmetry of the aperture: a square aperture produces a cross-shaped pattern with no rings, because the cancellation occurs along two axes rather than in circles.

\textit{``What determines the minimum size of a telescope---why can't we just make an infinitely large mirror?''}

The Airy disk tells you: the angular resolution is $\theta = 1.22\lambda/D$, so larger $D$ means finer resolution. But Feynman would note that the practical limits are not just diffraction. A mirror larger than about 10 meters faces engineering challenges: gravitational distortion (the mirror sags under its own weight), thermal expansion (temperature differences warp the surface), and atmospheric turbulence (the atmosphere blurs the image far worse than the diffraction limit for apertures above $\sim 20\,\text{cm}$). Modern large telescopes (Keck, VLT, the upcoming ELT) use segmented mirrors and adaptive optics to address these issues. The ultimate limit would be a space telescope (like Hubble or JWST) with a perfect mirror---there, the Airy disk is the true and final limit. JWST's 6.5 m primary mirror achieves $\theta_{\text{Airy}} \approx 0.02''$ at 2 $\mu$m, which is close to the diffraction limit.

\textit{``Why does the Airy disk contain 84\% of the light? Where does the rest go?''}

The remaining 16\% is distributed in the diffraction rings surrounding the central disk. The first ring contains 1.75\% of the total energy, the second ring 0.42\%, and so on, with each successive ring carrying less. The energy in the rings is ``wasted'' in the sense that it does not contribute to the image of the central source---it creates a faint halo that reduces contrast. For point sources (stars), the rings are usually too faint to see. For extended objects (planets, nebulae), the overlapping Airy patterns of adjacent points blur the image slightly. The fraction of energy in the central disk is $\eta = 1 - J_0^2(x_1) - J_1^2(x_1) \approx 0.838$, where $x_1 = 3.8317$ is the first zero of $J_1$. This is a fundamental property of the Bessel function and cannot be improved by better engineering.

\textit{``What happens to the Airy pattern if I use white light instead of monochromatic light?''}

Each wavelength produces its own Airy pattern with a different scale: $\theta(\lambda) = 1.22\lambda/D$, so red light ($\sim 700\,\text{nm}$) produces a pattern about 1.4 times larger than blue light ($\sim 500\,\text{nm}$). The central maxima overlap (the center is white), but the rings are colored: the first red ring is slightly outside the first blue ring, producing a rainbow-like halo. In practice, the overlapping patterns wash out the ring structure beyond the first ring or two, and the image is dominated by the central disk with a faint, blurred halo. For broadband imaging (like astronomical photography), a filter is often used to narrow the bandwidth and restore a well-defined Airy pattern. The chromatic blurring is one reason why high-resolution telescopes use narrow-band filters or spectrographs.

\textit{``Could a different aperture shape give better resolution than a circle?''}

Feynman would consider this carefully. A slit aperture produces a pattern that is very narrow in one direction ($\theta = \lambda/a$, where $a$ is the slit width) but very wide in the perpendicular direction. So in one direction, the slit is better than a circle of the same width; in the other direction, it is much worse. There is no aperture shape that beats a circle in \emph{all} directions simultaneously---this is a consequence of the uncertainty principle applied to the spatial frequency content of the aperture. However, for specific applications (like spectroscopy, where you need resolution in one direction), a slit is optimal. Apodization (tapering the aperture transmission from center to edge) can suppress the diffraction rings at the cost of broadening the central peak---useful when you need high contrast rather than raw resolution. The optimal aperture depends on what you are trying to measure.

\bigskip
\hrule
\bigskip
%=============================================================================
\section*{FAQ}

\textit{Q: Why 1.22 specifically---where does this number come from?}
A: The factor 1.22 arises from the first zero of the Bessel function $J_1(x)$, which occurs at $x = 3.8317$. Since $x = \pi D\sin\theta/\lambda$ and for small angles $\sin\theta \approx \theta$, the first dark ring is at $\theta = 3.8317/(\pi D/\lambda) = 1.22\lambda/D$. If the aperture were square instead of circular, the pattern would be $[\sin(x)/x]^2$ and the first zero would be at $\theta = \lambda/D$ (without the 1.22 factor). The 1.22 is purely a consequence of the circular geometry.

\textit{Q: Can we beat the Rayleigh criterion?}
A: Yes, with caveats. Super-resolution techniques (STED, PALM, STORM in microscopy) beat the diffraction limit by using fluorescent molecules that are switched on and off, localizing individual molecules to precision far better than $\lambda/2$. In astronomy, interferometric techniques (combining light from multiple telescopes) achieve the resolution of a telescope whose diameter equals the baseline between them. The Rayleigh criterion is a limit for \emph{incoherent} point sources viewed through a \emph{single} aperture; clever techniques can circumvent it by encoding additional information.

\textit{Q: What is the Airy disk of the human eye?}
A: The human pupil has a diameter of about 2--7 mm. At the peak sensitivity wavelength (550 nm), the Airy disk radius at the retina (focal length $\sim 17\,\text{mm}$) is $d = 2.44\lambda(f/D) \approx 5\,\mu\text{m}$ for a 3 mm pupil. This is comparable to the spacing of retinal cone cells ($\sim 2\,\mu\text{m}$), so the eye is approximately diffraction-limited. However, optical aberrations of the eye degrade the actual resolution to about $1'$ (one arcminute), which is larger than the diffraction limit for a 3 mm pupil.
\section*{Important Bibliography}
\begin{enumerate}
\item Born, M. and Wolf, E., \textit{Principles of Optics}, 7th ed. (Cambridge University Press, 1999), Chapter 8.
\item Hecht, E., \textit{Optics}, 5th ed. (Pearson, 2017), Chapter 10.
\item Airy, G.B., ``On the Diffraction of an Object-glass with Circular Aperture,'' \textit{Cambridge Philosophical Transactions} \textbf{5}, 283--291 (1835).
\item Goodman, J.W., \textit{Introduction to Fourier Optics}, 4th ed. (W.H. Freeman, 2017), Chapters 5--6.
\end{enumerate}


